\begin{statementbox}
\begin{center}
\begin{tikzpicture}[font=\small]
 \path[use as bounding box] (-4,-3) rectangle (4.5,3.2);
 \node at (0,2.85) {同一动点连到长轴两端点：一正一负，积恒定};
 \draw[->,gray] (-3.6,0)--(3.8,0) node[right] {$x$};
 \draw[->,gray] (0,-2.3)--(0,2.3) node[left] {$y$};
 \draw[blue!60!black,thick] (0,0) ellipse (3 and 2);
 \coordinate (A) at (-3,0); \coordinate (B) at (3,0); \coordinate (P) at (1.8,1.6);
 \draw[teal!70!black,very thick] (A)--(P);
 \draw[magenta!70!black,very thick] (B)--(P);
 \fill (A) circle (1.7pt) node[below left] {$A_1(-a,0)$};
 \fill (B) circle (1.7pt) node[below right] {$A_2(a,0)$};
 \fill (P) circle (1.7pt) node[above right] {$P(x,y)$};
 \node[teal!70!black] at (-1.4,1.1) {$k_1=\frac13$};
 \node[magenta!70!black,right] at (3,.9) {$k_2=-\frac43$};
 \node at (0,-2.65) {图取 $a=3,b=2$：$k_1k_2=-\frac49=-\frac{b^2}{a^2}$};
\end{tikzpicture}
\end{center}

\textbf{\textcolor{CStatement}{条件}}
\begin{description}[style=nextline,leftmargin=2.8em,labelsep=0.5em,itemsep=0.35em]
 \item[\textbf{\textcolor{CStatement}{条件一（长轴在 $x$ 轴）：}}]
  椭圆 $C:\dfrac{x^2}{a^2}+\dfrac{y^2}{b^2}=1\ (a>b>0)$，
  长轴端点 $A_1(-a,0),A_2(a,0)$；取 $P(x,y)\in C\setminus\{A_1,A_2\}$。
 \item[\textbf{\textcolor{CStatement}{条件二（斜率存在）：}}]
  由于 $P$ 不是两个长轴端点，必有 $-a<x<a$ 且 $y\ne0$，
  因而 $k_{PA_1},k_{PA_2}$ 都存在且一正一负。
\end{description}
\textbf{\textcolor{CStatement}{结论}}
\begin{description}[style=nextline,leftmargin=2.8em,labelsep=0.5em,itemsep=0.4em]
 \item[\textbf{\textcolor{CStatement}{结论一（椭圆第三定义）：}}]
  动点 $P$ 与\textbf{长轴两端点}连线的斜率乘积恒为负数：
  \[
   \boxed{k_{PA_1}\,k_{PA_2}=-\frac{b^2}{a^2}}
  \]
  其中 $b^2=a^2-c^2$，故还可写成 $\boxed{k_{PA_1}k_{PA_2}=e^2-1}$（$e=c/a$）。
 \item[\textbf{\textcolor{CStatement}{结论二（长轴在 $y$ 轴）：}}]
  若椭圆为 $\dfrac{y^2}{a^2}+\dfrac{x^2}{b^2}=1\ (a>b>0)$，
  长轴端点 $B_1(0,-a),B_2(0,a)$，且 $P$ 异于两端点，则
  \[
    \boxed{k_{PB_1}\,k_{PB_2}=-\frac{a^2}{b^2}}
  \]
  此时两条斜率均存在（$x\ne0$）。
 \item[\textbf{\textcolor{CStatement}{逆向认识（确定轨迹）：}}]
  固定 $A_1(-a,0),A_2(a,0)$（$a>0$），若动点 $P$ 使
  $k_{PA_1}k_{PA_2}=-\lambda$（$\lambda>0$），则轨迹位于椭圆
  \[
    \frac{x^2}{a^2}+\frac{y^2}{\lambda a^2}=1
  \]
  上并需剔除 $A_1,A_2$。只有 $0<\lambda<1$ 时，$A_1A_2$ 才是这条椭圆的\textbf{长轴}；
  $\lambda=1$ 为圆，$\lambda>1$ 时 $A_1A_2$ 是短轴。
\end{description}
\medskip
\textbf{\textcolor{CStatement}{圆与椭圆的双向坐标对应（理解结论的桥梁）：}}
取半径为 $a$ 的圆 $C_0:u^2+v^2=a^2$，与椭圆 $C$ 共用中心
$O(0,0)$ 和横轴端点 $A_1(-a,0),A_2(a,0)$。\textbf{特意选择保持横坐标不变}的纵向缩放：
\[
 \boxed{T:\ Q(u,v)\longmapsto P\left(u,\frac ba v\right)},
 \qquad
 \boxed{T^{-1}:\ P(x,y)\longmapsto Q\left(x,\frac ab y\right)}.
\]
这里 $0<b<a$；两式互逆，圆上每个点与椭圆上一个点一一对应（证明斜率积时须排除 $A_1,A_2$）。
它不是两条曲线之间\textbf{唯一可能}的对应，而是固定横轴、便于保留两端点并比较斜率的一种\textbf{特选映射}。
\textbf{不变量：}横坐标 $x=u$，横轴所有点（特别是 $A_1,O,A_2$），
以及对应点到横轴的垂足 $H=(u,0)$；归一化高度满足
$\dfrac{y}{b}=\dfrac{v}{a}$。长度和夹角通常不保持。

\medskip
\textbf{\textcolor{CStatement}{使用范围：}}
这是建立在椭圆标准方程和指定长轴端点上的\textbf{二级结论}（习惯称“第三定义”），
不可误作“椭圆上任意两定点均成立”。
\end{statementbox}
